\documentclass{beamer}
\usetheme{Madrid}
\usecolortheme{whale}
\usepackage{amsmath}
\usepackage{graphicx}
\usepackage{tikz}
\usepackage{pgfplots}
\usepackage{booktabs}
\usepackage{multirow}

\setbeamertemplate{footline}[frame number]{}
\setbeamertemplate{navigation symbols}{}
\titlegraphic{

\includegraphics[width=2cm]{SGAI/city-st-georges-logo.png}
}
\pgfplotsset{compat=1.18}


\title[DATE]{DATE: Derivative Alignment Training Extrapolation}
\author{E. Lopedoto, T. Weyde, and K. Salako}
\institute{City St George's, University of London}
\date{SGAI,Dec 2024}

\begin{document}

\begin{frame}
    \titlepage
\end{frame}

\begin{frame}{Problem: Neural Network Extrapolation}
    \begin{itemize}
        \item Neural Networks often struggle with extrapolation beyond training data
        \item We like ReLU
        \begin{enumerate}
            \item Natural linear extrapolation capability - unbounded
            \item Non-differentiable at x = 0 ("0-points")
            \item Cheap to compute
        \end{enumerate}
        \item Two key issues with ReLU networks:
        \begin{enumerate}
            \item Non-linearities outside training range produce no learning
            \item High variation in model derivative
        \end{enumerate}
    \end{itemize}
\end{frame}

\begin{frame}{Our contribution}
    \begin{figure}[t]
     \centering 
     \includegraphics[trim={1cm 1cm 1cm 1.9cm},clip, width=0.49\linewidth]{pics/SpagPred_pow_MSE__96355ae8-7134-11ef-8273-4b5d8eb7968f.eps}
     \includegraphics[trim={1cm 1cm 1cm 1.9cm},clip, width=0.49\linewidth]{pics/DerivPlot_pow_MSE__96355ae9-7134-11ef-8273-4b5d8eb7968f.eps}
    \caption{(a) Extrapolation with data trained in the range [-2,2]. (b) Interpolation range [-2,2] }
\end{figure}

    \begin{itemize}
        \item Solution to 1: Internal Range (IR) method
        \item Solution to 2: Derivative Loss (D1) method
    \end{itemize}
\end{frame}

\begin{frame}{ReLEx - Lopedoto et al., 2020}
    \begin{table}[t!]
     \centering
     \begin{tabular}{llcc} \hline
     Loss Name  & Definition & Complexity \\
     \hline \hline
     CP - Centripetal  & $\sum_{h=1}^H \sum_{k=1}^K (x_{h0} - x_{k})^2$ & $O(H K)$ \\
     \hfill
     \\
     MR - Mutually Repellent & $\sum_{h=1}^H \sum_{j>i}^H \frac{1}{(x_{h0} - x_{j0})^2+ \varepsilon}$ & $O(H^2)$ \\ 
     \hfill
     \\
     WO - Weight Orientation & $- \sum_{h=1}^H {ReLU((x_{h0} - \bar{x})\cdot w_h + \varepsilon)}$ & $O(H)$ \\
     \hfill
     \\
     WS - Weight Sign $ & $({\sum_{h=1}^Hw_h})^2$ & $O(H)$ 
     \\
     \hline
    \end{tabular}
    
    \caption{The loss terms used in ReLEx improve extrapolation. 
    $K$ is the number of data points.
    $H$ is the number of hidden neurons. 
    $x_{h0}$ is the \emph{0-point} of hidden neuron $h$ (dependent only on the network weights $w_h$ and $b_h$).
    $\bar{x}$ is the average of the training data. 
    The complexity refers to a single calculation of the loss value. 
    }
     
    \end{table}
\end{frame}

\begin{frame}{The DATE Method}
    Two new regularization terms:
    \begin{enumerate}
        \item IR-loss ($L_{IR}$):
        \begin{itemize}
            \item Keeps 0-points inside training data range
            \item More efficient than ReLEx approach
        \end{itemize}
        \item D1-loss ($L_{D1}$):
        \begin{itemize}
            \item Penalizes deviation of model derivative from target derivative
            \item Estimated from training data points
        \end{itemize}
        
    \end{enumerate}
\end{frame}

\begin{frame}{IR-Loss Implementation}
    \begin{itemize}
        \item Purpose: Keep 0-points within training range
        \item Formula:
        \[ L_{IR} = \sum_{k=1}^K (ReLU(-x_{k0} + L - \tau) + ReLU(x_{k0} - R + \tau))^2 \]
        where:
        \begin{itemize}
            \item $L$, $R$: Left and right bounds of training range
            \item $\tau$: Margin parameter
            \item $x_{k0}$: 0-point of neuron k
        \end{itemize}
        \item Complexity: O(K) vs O(H²) in ReLEx
    \end{itemize}
\end{frame}

\begin{frame}{D1-Loss Implementation}
    \begin{itemize}
        \item Derivative estimation:
        \[ y'_k = \frac{y_{k+1} - y_k}{x_{k+1} - x_k} \]
        \item Model derivative:
        \[ \hat{f}'(x) = \frac{\hat{f}(x + \epsilon) - \hat{f}(x - \epsilon)}{2\epsilon} \]
        \item D1-loss:
        \[ L_{D1} = \sum_{(x',y')\in XY'} (y' - \hat{f}'(x'))^2 \]
    \end{itemize}
\end{frame}




%\begin{frame}{Background: ReLU Networks}
%    \begin{itemize}
%        \item Properties:
%        \begin{itemize}
%            \item Natural linear extrapolation capability
%            \item Non-differentiable at x = 0 ("0-points")
%            \item "Dying ReLU" problem when inputs are always negative
%        \end{itemize}
%        \item Simple feedforward network:
%        \[ f(x,\beta) = b_o + \sum_{h=1}^H w_{ho} \cdot ReLU(w_h x + b_h) \]
%    \end{itemize}
%\end{frame}

\begin{frame}{Experimental Setup}
    \begin{itemize}
        \item Network: Single hidden layer, 10 neurons
        \item ReLU activation: $ReLU(x) = max(0,x)$

        \item Training data:
        \begin{itemize}
            \item Range: [-2, 2]
            \item 200 equidistant points
            \item Test functions: linear, absolute, quadratic, sigmoid, sine
        \end{itemize}
        \item Hyperparameters:
        \begin{itemize}
            \item Learning rate: 0.007
            \item Momentum: 0.9
            \item 700 epochs, batch size 20
        \end{itemize}
    \end{itemize}
\end{frame}


\begin{frame}{Evaluation Metrics}
    \begin{itemize}
        \item Interpolation metrics:
        \begin{itemize}
            \item $MSE_{int}$: Error in training range
            \item $MSE_{D1}$: Derivative error in training range
        \end{itemize}
        \item Extrapolation metrics:
        \begin{itemize}
            \item $Tan_{ex}$: Error vs tangent lines
            \item $D1_{ex}$: Error vs last two points slope
            \item $WLR_{ex}$: Error vs weighted linear regression
            
        \end{itemize}
        \item Consistency: $\sigma_S$ (slope standard deviation)
    \end{itemize}
\end{frame}


\begin{frame}{DATE Configurations}
    Three variants tested:
    \begin{enumerate}
        \item $DATE_1$: Full ReLEx + $L_{D1}$
        \item $DATE_2$: $L_{WO} + L_{WS} + L_{IR} + L_{D1}$
        \item $DATE_3$: $L_{WS} + L_{IR} + L_{D1}$
    \end{enumerate}
    Key differences:
    \begin{itemize}
        \item $DATE_1$: Most complex but stable
        \item $DATE_2$: Balanced performance
        \item $DATE_3$: Simplest configuration
    \end{itemize}
\end{frame}


\begin{frame}{Results: Method Comparison - Pow function}
   \begin{table}
     \centering
     \begin{tabular}{l|l|rrr|Rrr}
     $f(x)$ & Loss & $\sigma_{S}$ & 
     $MSE_{int}$ & $MSE_{D1}$ & $Tan_{ex}$ & $WLR_{ex}$ & $D1_{ex}$ \\
     \hline 
     
    \multirow{5}{0.8cm}{$pow$} 
     & $MSE$ &  0.479 & 0.081 & 0.808 & 8.974 & 5.317 & 8.958 \\
     & $ReLEx$ & 0.014 & 0.036 & 0.673 & 2.970 & 0.244 & 2.955 \\
     & $DATE_1$ & \textbf{0.017} & 0.038 & 0.511 & 2.376 & \textbf{0.460} & 2.362 \\
     & $DATE_2$ &  0.021 & \textbf{0.007} & 0.108 & \textbf{0.141} & 3.463 & \textbf{0.138} \\
     & $DATE_3$ & 0.035 & 0.011 &\textbf{ 0.097} & 0.206 & 3.272 & 0.202 \\
     \hline
     \hline
     \end{tabular}
    \end{table}

    \begin{figure}
     \centering 

     \includegraphics[trim={0.5cm 0cm 1cm 2cm},clip, width=0.3\linewidth]{pics/DerivPlot_pow_ReLEx__4eeb8531-7135-11ef-8273-4b5d8eb7968f.eps}
     \includegraphics[trim={0.5cm 0cm 1cm 2cm},clip, width=0.3\linewidth]{pics/DerivPlot_pow_ReLEx__b0ce38b1-7135-11ef-8273-4b5d8eb7968f.eps}
     \includegraphics[trim={0.5cm 0cm 1cm 2cm},clip, width=0.3\linewidth]{pics/DerivPlot_pow_ReLEx__15c78ebb-7136-11ef-8273-4b5d8eb7968f.eps}
    
    
     \begin{tabular}{p{0.3\linewidth} p{0.3\linewidth} p{0.3\linewidth}}
    $DATE_1$ & $DATE_2$ & $DATE_3$
     
     \end{tabular} 
    \end{figure}
\end{frame}


\begin{frame}{Conclusion}
    \begin{itemize}
        \item DATE improves extrapolation while maintaining good interpolation
        \item Key advantages:
        \begin{itemize}
            \item Better derivative alignment
            \item More efficient than ReLEx
            \item Multiple configuration options
        \end{itemize}
    
    \end{itemize}
\end{frame}


\begin{frame}{Improvements \& Future Work}
    \begin{itemize}
        \item Current limitations:
        \begin{itemize}
            \item Only univariate regression
            \item Single hidden layer
            \item Assumes linear extrapolation is desired
            \item Sensitivity to noise at margins
        \end{itemize}
        \item Future work:
        \begin{itemize}
            \item Extension to multiple dimensions
            \item Different network architectures
            \item Higher-order derivatives
            \item Working with Smoothing
        \end{itemize}
    \end{itemize}
\end{frame}




\end{document}